\section{The Decision Procedure}
\label{sec:decision}

The two cost laws of \S\ref{sec:cost_laws} produce two crossovers, and those
crossovers delimit the regime in which block-structured decoding wins on
per-event latency. This section turns that derivation into an executable
artifact: a procedure that takes four deployment inputs and returns which
architecture to serve. It is implemented, tested, and runnable from a clean
checkout with no GPU.

\paragraph{Inputs.}
The procedure takes $(L, \lambda, B, \tau)$: expected output length $L$
(tokens), arrival rate $\lambda$ (events/s), batching availability
$B\in\{0,1\}$, and Field-F1 tolerance $\tau$ (the F1 a caller will trade for
latency). Every constant the procedure uses traces to a committed artifact and
a \texttt{NUMBERS.md} section; none is hardcoded in the result.

\paragraph{Cost laws in force.}
From \S\ref{sec:cost_laws},
\[
T_{\textsc{ar}}(L)=a+bL,\qquad a=\arBase,\ b=\arSlope\ \text{ms/tok},
\]
\[
T_{\textsc{dg}}(L)=c+s\!\left(\left\lceil L/C\right\rceil-1\right),\qquad
c=\dgBase,\ s=\dgStep,\ C=\canvasLength.
\]
The DG base is prefill-dominated ($\dgPrefillPct\%$ prefill $+$ fixed,
$\dgDenoisePct\%$ denoising at $K{=}8$), so the staircase is \emph{structural}:
a property of the canvas boundary, not of the diffusion process.

\paragraph{Crossovers.}
Solving $T_{\textsc{ar}}(L^\star)=T_{\textsc{dg}}(L^\star)$ at the first
canvas ($\lceil L^\star/C\rceil=1$, so $T_{\textsc{dg}}=c$) gives
\[
L^\star=\frac{c-a}{b}=\frac{\dgBase-\arBase}{\arSlope}\approx \crossoverL\ \text{tokens}.
\]
DG's serial throughput is the bar AR batching must clear:
\[
\lambda^\star=\frac{1000}{c}=\frac{1000}{\dgBase}\approx \crossoverLambda\ \text{ev/s}.
\]
Both are $K$-dependent: at $K{=}32$ ($c=\dgOldBase$), $L^\star\approx 65$ and
$\lambda^\star\approx 2.4$. The procedure reports the $K{=}8$ values as
default and reproduces both from the law constants, not from hardcoded
answers.

\paragraph{The procedure.}
Five guards fire in order; the first that fires wins
(Algorithm~\ref{alg:decide}).

\begin{algorithm}[ht]
\caption{Decide which architecture to serve.}
\label{alg:decide}
\begin{algorithmic}[1]
\State $T_{\textsc{ar}}\gets a+bL$;\ \ $T_{\textsc{dg}}\gets c+s(\lceil L/C\rceil-1)$;\ \ $r\gets T_{\textsc{ar}}/T_{\textsc{dg}}$
\If{$L < L_{\mathrm{nat}}^{\textsc{dg}}$} \Comment{guard 1: truncation}
  \State \Return \textsc{ar} \Comment{DG output truncated below natural length; quality collapses}
\ElsIf{$B \wedge \lambda > \lambda^\star$} \Comment{guard 2: throughput}
  \State \Return \textsc{ar} \Comment{DG serial throughput saturated; AR batching wins}
\ElsIf{$L < L^\star$} \Comment{guard 3: length}
  \State \Return \textsc{ar} \Comment{below the latency crossover}
\ElsIf{$\tau < |\Delta F_1|$} \Comment{guard 4: quality}
  \State \Return \textsc{ar} \Comment{quality price exceeds tolerance}
\Else \Comment{guard 5: default}
  \State \Return \textsc{dg} with expected speedup $r$
\EndIf
\end{algorithmic}
\end{algorithm}

Guard~3 is structurally subsumed by guard~1 in the current configuration
($L_{\mathrm{nat}}^{\textsc{dg}}=176 > L^\star=\crossoverL$): any output short
enough to lie below the latency crossover is already below DG's natural
length, so truncation fires first. Both recommend \textsc{ar}, so the
recommendation is unchanged; the length guard remains because it is the guard
that fires once a shorter-canvas DG variant is configured. The quality guard's
threshold is metric-dependent: on the contaminated 12-field metric the price
is $|\Delta F_1|=0.039$, but on the 9-field metric (\S\ref{sec:exchange_rate}) the
price is $\approx 0$ ($\arThreeBFnine$ vs $\dgNineF$), so the guard is inert
and DG pays no quality tax at equal quality.

\paragraph{Worked examples.}
Table~\ref{tab:worked} walks six scenarios through the procedure. The ratio
column is $r=T_{\textsc{ar}}/T_{\textsc{dg}}$---the per-event speedup DG
\emph{would} deliver---reported even when a non-latency guard selects
\textsc{ar}, so the opportunity cost of each guard is visible.

\begin{table}[ht]
\centering
\caption{Six worked examples through the decision procedure.}
\label{tab:worked}
\small
\begin{tabular}{lrrrrlll}
\toprule
Scenario & $L$ & $\lambda$ & $B$ & $\tau$ & Arm & $T_{\textsc{ar}}/T_{\textsc{dg}}$ & Rule \\
\midrule
Interactive single-record extraction      & 176  & 0.5 & no & 0.05 & \textbf{DG} & $2.94\times$  & default \\
Long stack-trace summarisation, batch-1   & 2048 & 1.0 & no & 0.05 & \textbf{DG} & $10.6\times$ & default \\
High-volume log pipeline                   & 256  & 64  & yes& 0.05 & \textbf{AR} & $4.27\times$ & throughput \\
Short field-completion                     & 32   & 0.5 & no & 0.10 & \textbf{AR} & $0.55\times$ & truncation \\
Tight output budget, unknown record size   & 128  & 0.5 & no & 0.05 & \textbf{AR} & $2.14\times$ & truncation \\
Quality-critical extraction                & 512  & 1.0 & no & 0.01 & \textbf{AR} & $6.48\times$ & quality \\
\bottomrule
\end{tabular}
\end{table}

Rows~4--5 carry the key subtlety: in row~5 DG would be $2.14\times$ faster
per-event, yet \textsc{ar} is chosen, because at $L<176$ DG is truncated below
its natural output length and its quality collapses. The guard protects
\emph{quality}, not latency; the ratio column exposes the latency an operator
gives up to obtain valid output.

\paragraph{Implementation.}
The procedure lives in \texttt{models/decision/procedure.py}, with a CLI:
\begin{center}
\texttt{python -m models.decision -{}-L 512 -{}-lambda 1.0 -{}-no-batching -{}-f1-tolerance 0.05}
\end{center}
$T_{\textsc{dg}}$ is \emph{injectable}: the one-variable law $T_{\textsc{dg}}(L)$
ships now, and the two-variable law $T_{\textsc{dg}}(L_{\mathrm{in}},L_{\mathrm{out}})$
swaps in as the only change when the input-length sweep
(\S\ref{sec:cost_laws}, two-variable scope) is in force---in two-variable mode
the procedure emits a warning that its prefill slope is pending calibration.
The self-test (\texttt{test\_procedure.py}) reproduces $L^\star$ and
$\lambda^\star$ from the law constants, fires each guard in isolation, checks
all six worked examples, and probes step-range robustness: \textbf{11/11 tests
pass} (claim~C15).

\paragraph{Validity domain.}
The constants are calibrated on one \gpuModel{} under vLLM \vllmVersion{},
BF16, serial diffusion (\texttt{-{}-max-num-seqs 1}), for structured extraction
over BGL with a $\sim$1047-token few-shot prompt. The procedure must not be
applied outside this domain---in particular to batch-served diffusion,
canvases other than \canvasLength, materially different input lengths, other
GPUs/runtimes, or 12-field F1 tolerances (use the 9-field metric).
